\scanpage{143}
\mspage{1}
1. \uline{Cofinalité} \uline{cohomologique}

Soit
\[
 \alpha:C'\longrightarrow C
\]
un foncteur de (petites) Catégories, $M$ une catégorie ayant des $\varprojlim$ \unsure{pertinentes}, $F$ un préfaisceau\inserted{\unsure{donc contravariant ici}} sur $C$ à val. dans $M$. On a \unsure{un hom.|une trans.} can. dans $M$
\[
 \varprojlim_C F\longrightarrow\varprojlim_{C'}\underbrace{\alpha^*(F)}_{F\circ\alpha}
 \tag*{(1)}
\]
et \uline{\unsure{Ch.|Cher.}} \deleted{\unsure{cond}} des conditions sur $\alpha$ pour que ce soit tjrs un iso. On est ramené aussitôt au cas de $M=(\mathrm{Ens})$.
% slanted
\margin{On dit alors que $\alpha$ est \uuline{cofinal}, cf SGA 4, \unsure{4 I|I} \S~8.1. ou 3.3. p. 45}
\deleted{On dit \unsure{parfois que} $\alpha$ est \unsure{cofinal} si pour tt \unsure{$x$}} \deleted{On a \unsure{$C$}}

Disons qu'une catégorie $\unsure{\mathbb I|I}$ est \inserted{\uline{fbt}} \uline{filtr. à g} si \inserted{$I\ne\varnothing$ et} si $\forall\,\unsure{i'|x'},\unsure{i''|x''}\in\operatorname{Ob}I$, $\exists\,i\in\operatorname{Ob}I$ et des flèches
\diagram{\begin{tikzcd}[ampersand replacement=\&,row sep=small,column sep=small]
\&i'\\i\arrow[ur]\arrow[dr]\&\\ \&i''
\end{tikzcd}}
Cela signifie aussi que $\operatorname{Ob}I$ en tant qu'ens ordonné est filtrant à gauche.

\uline{Proposition} Pour que (1) soit toujours un iso, \inserted{($\alpha$ ``anodin pour $\varprojlim$'')} il suffit que pour tout $x_0\in\operatorname{Ob} C$, la catégorie ${}_{x_0\backslash}\overset{\vee}{C'}^{\unsure{x_0|y_x}}$ des couples $(x',u)$, avec $x'\in\operatorname{Ob}C'$, $u:\alpha(x')\longrightarrow x_0$, soit fbt. filt. à g.
% slanted
\margin{\unsure{cf condition bien|amélioration bien} \unsure{moins que|moins exigeante} p. 5. \quad cf 3.2 p. 43--44 pour \unsure{caractérisation|condition} de cofinalité}

Ceci n'implique pas la conclusion plus forte
\[
 F\xrightarrow[\mathrm{adj.}]{\sim}\alpha_*\alpha^*F
\]
que le morphisme d'adj. canonique soit un iso. Pour qu'il en soit ainsi, il \deleted{faut} \inserted{suffit} que $\forall x_0\in C$, le foncteur \inserted{$\alpha_{/x_0}:$} $C'_{/x_0}\unsure{-|\to}C_{/x_0}$ induit par $\alpha$ soit ``anodin pour les $\varprojlim$'', ce qui signifie aussi que pour toute flèche $u:x_1\longrightarrow x_0$ dans $C$, la catégorie $\underset{u}{\underset{\vee}{C'}}$ formée des triplets $(x',\zeta,\eta)$ avec $x'\in\operatorname{Ob}C'$ et $\zeta:x_1-\alpha(x')$, $\eta:\alpha(x')\longrightarrow x_0$ rendant commutatif
% archive sheet number: 143
